\ifdefined\gkver\else\def\gkver{student}\fi
\documentclass[\gkver]{gaokaozhenti}
\begin{document}

\chapter{2024年山东}

\begin{enumerate}
\item
%% number: 1
%% paperName: 2024年山东高考第 1 题 3 分
%% typeId: 1
%% type: 单选题
%% chapter: 原子和原子核
%% point: 半衰期
%% method: 无
%% score: 3
%% degree: 750
%% duplicateId: 0
%% body:
2024 年是中国航天大年，神舟十八号、嫦娥六号等已陆续飞天，部分航天器装载了具有抗干扰性强的核电池。已知 \ce{^{90}_{38}Sr} 衰变为 \ce{^{90}_{39}Y} 的半衰期约为 29 年；\ce{^{238}_{94}Pu} 衰变为 \ce{^{234}_{92}U} 的半衰期约 87 年。现用相同数目的 \ce{^{90}_{38}Sr} 和 \ce{^{238}_{94}Pu} 各做一块核电池，下列说法正确的是\xzanswer{D}

\fourchoices[answer=D]
{\ce{^{90}_{38}Sr} 衰变为 \ce{^{90}_{39}Y} 时产生 $\alpha$ 粒子}
{\ce{^{238}_{94}Pu} 衰变为 \ce{^{234}_{92}U} 时产生 $\beta$ 粒子}
{50 年后，剩余的 \ce{^{90}_{38}Sr} 数目大于 \ce{^{238}_{94}Pu} 的数目}
{87 年后，剩余的 \ce{^{90}_{38}Sr} 数目小于 \ce{^{238}_{94}Pu} 的数目}

%% answer: D
%% memo:
\memoanswer{
A．根据质量数守恒和电荷数守恒可知 \ce{^{90}_{38}Sr} 衰变为 \ce{^{90}_{39}Y} 时产生电子，即 $\beta$ 粒子，故 A 错误；

B．根据质量数守恒和电荷数守恒可知 \ce{^{238}_{94}Pu} 衰变为 \ce{^{234}_{92}U} 时产生 \ce{^{4}_{2}He}，即 $\alpha$ 粒子，故 B 错误；

CD．根据题意可知 \ce{^{238}_{94}Pu} 的半衰期大于 \ce{^{90}_{38}Sr} 的半衰期，现用相同数目的 \ce{^{90}_{38}Sr} 和 \ce{^{238}_{94}Pu} 各做一块核电池，经过相同的时间，\ce{^{90}_{38}Sr} 经过的半衰期的次数多，所以 \ce{^{90}_{38}Sr} 数目小于 \ce{^{238}_{94}Pu} 的数目，故 D 正确，C 错误。

故选 D。
}

\item
%% number: 2
%% paperName: 2024年山东高考第 2 题 3 分
%% typeId: 1
%% type: 单选题
%% chapter: 力物体的平衡
%% point: 三力平衡
%% method: 无
%% score: 3
%% degree: 850
%% duplicateId: 0
%% body:
如图所示，国产人形机器人“天工”能平稳通过斜坡。若它可以在倾角不大于 $30^{\circ}$ 的斜坡上稳定地站立和行走，且最大静摩擦力等于滑动摩擦力，则它的脚和斜面间的动摩擦因数不能小于\xzanswer{B}

\onepicture[w=4.20cm, align=right]{figs/02.jpg}

\fourchoices[answer=B]
{$\frac{1}{2}$}
{$\frac{\sqrt{3}}{3}$}
{$\frac{\sqrt{2}}{2}$}
{$\frac{\sqrt{3}}{2}$}

%% answer: B
%% memo:
\memoanswer{
根据题意可知机器人“天工”它可以在倾角不大于 $30^{\circ}$ 的斜坡上稳定地站立和行走，对“天工”分析有

$mg\sin30^{\circ} \le \mu mg\cos30^{\circ}$

可得 $\mu \ge \frac{\sqrt{3}}{3}$

故选 B。
}

\item
%% number: 3
%% paperName: 2024年山东高考第 3 题 3 分
%% typeId: 1
%% type: 单选题
%% chapter: 直线运动
%% point: 匀变速直线运动
%% method: 无
%% score: 3
%% degree: 750
%% duplicateId: 0
%% body:
如图所示，固定的光滑斜面上有一木板，其下端与斜面上 A 点距离为 $L$。木板由静止释放，若木板长度 $L$，通过 A 点的时间间隔为 $\Delta t_{1}$；若木板长度为 $2L$，通过 A 点的时间间隔为 $\Delta t_{2}$。$\Delta t_{2}:\Delta t_{1}$ 为\xzanswer{A}

\onepicture[w=4.50cm, align=right]{figs/03.png}

\fourchoices[answer=A]
{($\sqrt{3}-1$)$:$($\sqrt{2}-1$)}
{($\sqrt{3}-\sqrt{2}$)$:$($\sqrt{2}-1$)}
{($\sqrt{3}+1$)$:$($\sqrt{2}+1$)}
{($\sqrt{3}+\sqrt{2}$)$:$($\sqrt{2}+1$)}

%% answer: A
%% memo:
\memoanswer{
木板在斜面上运动时，木板的加速度不变，设加速度为 $a$，木板从静止释放到下端到达 A 点的过程，根据运动学公式有 $L = \frac{1}{2}at_{0}^{2}$；

木板从静止释放到上端到达 A 点的过程，当木板长度为 $L$ 时，有 $2L = \frac{1}{2}at_{1}^{2}$；

当木板长度为 $2L$ 时，有 $3L = \frac{1}{2}at_{2}^{2}$；

又 $\Delta t_{1} = t_{1} - t_{0}$，$\Delta t_{2} = t_{2} - t_{0}$

联立解得 $\Delta t_{2}:\Delta t_{1} = (\sqrt{3}-1):(\sqrt{2}-1)$

故选 A。
}

\item
%% number: 4
%% paperName: 2024年山东高考第 4 题 3 分
%% typeId: 1
%% type: 单选题
%% chapter: 光学
%% point: 光的干涉
%% method: 无
%% score: 3
%% degree: 850
%% duplicateId: 0
%% body:
检测球形滚珠直径是否合格的装置如图甲所示，将标准滚珠 a 与待测滚珠 b、c 放置在两块平板玻璃之间，用单色平行光垂直照射平板玻璃，形成如图乙所示的干涉条纹。若待测滚珠与标准滚珠的直径相等为合格，下列说法正确的是\xzanswer{C}

\onepicture[w=7.50cm]{figs/04.png}

\fourchoices[answer=C]
{滚珠 b、c 均合格}
{滚珠 b、c 均不合格}
{滚珠 b 合格，滚珠 c 不合格}
{滚珠 b 不合格，滚珠 c 合格}

%% answer: C
%% memo:
\memoanswer{
单色平行光垂直照射平板玻璃，上、下玻璃上表面的反射光在上玻璃上表面发生干涉，形成干涉条纹，光程差为两块玻璃距离的两倍，根据光的干涉知识可知，同一条干涉条纹位置处光的波程差相等，即滚珠 $a$ 的直径与滚珠 $b$ 的直径相等，即滚珠 $b$ 合格，不同的干涉条纹位置处光的波程差不同，则滚珠 $a$ 的直径与滚珠 $c$ 的直径不相等，即滚珠 $c$ 不合格。

故选 C。
}

\item
%% number: 5
%% paperName: 2024年山东高考第 5 题 3 分
%% typeId: 1
%% type: 单选题
%% chapter: 曲线运动
%% point: 开普勒定律
%% method: 无
%% score: 3
%% degree: 550
%% duplicateId: 0
%% body:
“鹊桥二号”中继星环绕月球运行，其 24 小时椭圆轨道的半长轴为 $a$。已知地球同步卫星的轨道半径为 $r$，则月球与地球质量之比可表示为\xzanswer{D}

\fourchoices[answer=D]
{$\sqrt{\frac{r^{3}}{a^{3}}}$}
{$\sqrt{\frac{a^{3}}{r^{3}}}$}
{$\frac{r^{3}}{a^{3}}$}
{$\frac{a^{3}}{r^{3}}$}

%% answer: D
%% memo:
\memoanswer{
“鹊桥二号”中继星在 24 小时椭圆轨道运行时，根据开普勒第三定律 $\frac{a^{3}}{T^{2}} = k$

同理，对地球的同步卫星根据开普勒第三定律 $\frac{r^{3}}{T'^{2}} = k'$

又开普勒常量与中心天体的质量成正比，所以 $\frac{M_{\text{月}}}{M_{\text{地}}} = \frac{k}{k'}$

联立可得 $\frac{M_{\text{月}}}{M_{\text{地}}} = \frac{a^{3}}{r^{3}}$

故选 D。
}

\item
%% number: 6
%% paperName: 2024年山东高考第 6 题 3 分
%% typeId: 1
%% type: 单选题
%% chapter: 气体性质
%% point: 热力学第一定律
%% method: 无
%% score: 3
%% degree: 550
%% duplicateId: 0
%% body:
一定质量理想气体经历如图所示的循环过程，$a \to b$ 过程是等压过程，$b \to c$ 过程中气体与外界无热量交换，$c \to a$ 过程是等温过程。下列说法正确的是\xzanswer{C}

\onepicture[w=4.60cm, align=right]{figs/06.png}

\fourchoices[answer=C]
{$a \to b$ 过程，气体从外界吸收的热量全部用于对外做功}
{$b \to c$ 过程，气体对外做功，内能增加}
{$a \to b \to c$ 过程，气体从外界吸收的热量全部用于对外做功}
{$a \to b$ 过程，气体从外界吸收的热量等于 $c \to a$ 过程放出的热量}

%% answer: C
%% memo:
\memoanswer{
A．$a \to b$ 过程压强不变，是等压变化且体积增大，气体对外做功 $W < 0$，由盖-吕萨克定律可知 $T_{b} > T_{a}$

即内能增大，$\Delta U_{ab} > 0$，根据热力学第一定律 $\Delta U = Q + W$ 可知 $a \to b$ 过程，气体从外界吸收的热量一部分用于对外做功，另一部分用于增加内能，A 错误；

B．方法一：$b \to c$ 过程中气体与外界无热量交换，即 $Q_{bc} = 0$

又由气体体积增大可知 $W_{bc} < 0$，由热力学第一定律 $\Delta U = Q + W$ 可知气体内能减少。

方法二：$c \to a$ 过程为等温过程，所以 $T_{c} = T_{a}$。

结合 $T_{b} > T_{a}$ 分析可知 $T_{b} > T_{c}$。

所以 $b$ 到 $c$ 过程气体的内能减少。故 B 错误；

C．$c \to a$ 过程为等温过程，可知 $T_{c} = T_{a}$，$\Delta U_{ac} = 0$。

根据热力学第一定律可知 $a \to b \to c$ 过程，气体从外界吸收的热量全部用于对外做功，C 正确；

D．根据热力学第一定律结合上述解析可知：$a \to b \to c \to a$ 一整个热力学循环过程 $\Delta U = 0$，整个过程气体对外做功，因此热力学第一定律可得

$\Delta U = Q_{ab} - Q_{ca} - W = 0$

故 $a \to b$ 过程气体从外界吸收的热量 $Q_{ab}$ 不等于 $c \to a$ 过程放出的热量 $- Q_{ca}$，D 错误。

故选 C。
}

\item
%% number: 7
%% paperName: 2024年山东高考第 7 题 3 分
%% typeId: 1
%% type: 单选题
%% chapter: 功和能
%% point: 功能原理
%% method: 无
%% score: 3
%% degree: 550
%% duplicateId: 0
%% body:
如图所示，质量均为 $m$ 的甲、乙两同学，分别坐在水平放置的轻木板上，木板通过一根原长为 $l$ 的轻质弹性绳连接，连接点等高且间距为 $d$（$d < l$）。两木板与地面间动摩擦因数均为 $\mu$，弹性绳劲度系数为 $k$，被拉伸时弹性势能 $E = \frac{1}{2}kx^{2}$（$x$ 为绳的伸长量）。现用水平力 $F$ 缓慢拉动乙所坐木板，直至甲所坐木板刚要离开原位置，此过程中两人与所坐木板保持相对静止，$k$ 保持不变，最大静摩擦力等于滑动摩擦力，重力加速度大小为 $g$，则 $F$ 所做的功等于\xzanswer{B}

\onepicture[w=8.00cm, align=right]{figs/07.png}

\fourchoices[answer=B]
{$\frac{(\mu mg)^{2}}{2k} + \mu mg(l - d)$}
{$\frac{3(\mu mg)^{2}}{2k} + \mu mg(l - d)$}
{$\frac{3(\mu mg)^{2}}{2k} + 2\mu mg(l - d)$}
{$\frac{(\mu mg)^{2}}{2k} + 2\mu mg(l - d)$}

%% answer: B
%% memo:
\memoanswer{
当甲所坐木板刚要离开原位置时，对甲及其所坐木板整体有 $kx_{0} = \mu mg$

解得弹性绳的伸长量 $x_{0}=\frac{\mu mg}{k}$

则此时弹性绳的弹性势能为

$E_{0}=\frac{1}{2}kx_{0}^{2}=\frac{\mu^{2}m^{2}g^{2}}{2k}$

从开始拉动乙所坐木板到甲所坐木板刚要离开原位置的过程，乙所坐木板的位移为 $x_{1}=x_{0}+l-d$

则由功能关系可知该过程 $F$ 所做的功

$W = E_{0} + \mu mgx_{1} = \frac{3(\mu mg)^{2}}{2k} + \mu mg(l - d)$

故选 B。
}

\item
%% number: 8
%% paperName: 2024年山东高考第 8 题 3 分
%% typeId: 1
%% type: 单选题
%% chapter: 电磁感应
%% point: 切割磁感线电动势
%% method: 无
%% score: 3
%% degree: 550
%% duplicateId: 0
%% body:
如图甲所示，在 $- d \le x \le d$，$- d \le y \le d$ 的区域中存在垂直 $Oxy$ 平面向里、磁感应强度大小为 $B$ 的匀强磁场（用阴影表示磁场的区域），边长为 $2d$ 的正方形线圈与磁场边界重合。线圈以 $y$ 轴为转轴匀速转动时，线圈中产生的交变电动势如图乙所示。若仅磁场的区域发生了变化，线圈中产生的电动势变为图丙所示实线部分，则变化后磁场的区域可能为\xzanswer{C}

\onepicture[w=8.50cm]{figs/08a.png}

\fourchoices[answer=C, ispicture=true, h=2.7cm]
{figs/08b1.png}
{figs/08b2.png}
{figs/08b3.png}
{figs/08b4.png}

%% answer: C
%% memo:
\memoanswer{
根据题意可知，磁场区域变化前线圈产生的感应电动势为 $e = E\sin\omega t$

由题图丙可知，磁场区域变化后，当 $E\sin\omega t = \frac{\sqrt{3}}{2}E$ 时，线圈的侧边开始切割磁感线，即当线圈旋转 $\frac{\pi}{3}$ 时开始切割磁感线，由几何关系可知磁场区域平行于 $x$ 轴的边长变为

$d' = 2d\cos\frac{\pi}{3} = d$

C 正确。

故选 C。
}

\item
%% number: 9
%% paperName: 2024年山东高考第 9 题 4 分
%% typeId: 2
%% type: 多选题
%% chapter: 机械波
%% point: 波的图象
%% method: 无
%% score: 4
%% degree: 550
%% duplicateId: 0
%% body:
甲、乙两列简谐横波在同一均匀介质中沿 $x$ 轴相向传播，波速均为 $2 \Ums$。$t = 0$ 时刻二者在 $x = 2 \Um$ 处相遇，波形图如图所示。关于平衡位置在 $x = 2 \Um$ 处的质点 P，下列说法正确的是\xzanswer{BC}

\onepicture[w=6.00cm, align=right]{figs/09.png}

\fourchoices[answer=BC]
{$t = 0.5 \Us$ 时，P 偏离平衡位置的位移为 0}
{$t = 0.5 \Us$ 时，P 偏离平衡位置的位移为 $- 2 \Ucm$}
{$t = 1.0 \Us$ 时，P 向 $y$ 轴正方向运动}
{$t = 1.0 \Us$ 时，P 向 $y$ 轴负方向运动}

%% answer: BC
%% memo:
\memoanswer{
AB．在 $0.5 \Us$ 内，甲、乙两列波传播的距离均为

$\Delta x = v\Delta t = 2 \times 0.5 \Um = 1 \Um$

根据波形平移法可知，$t = 0.5 \Us$ 时，$x = 1 \Um$ 处甲波的波谷刚好传到 P 处，$x = 3 \Um$ 处乙波的平衡位置振动刚好传到 P 处，根据叠加原理可知，$t = 0.5 \Us$ 时，P 偏离平衡位置的位移为 $- 2 \Ucm$，故 A 错误，B 正确；

CD．在 $1.0 \Us$ 内，甲、乙两列波传播的距离均为

$\Delta x' = v\Delta t' = 2 \times 1.0 \Um = 2 \Um$

根据波形平移法可知，$t = 1.0 \Us$ 时，$x = 0$ 处甲波的平衡位置振动刚好传到 P 处，$x = 4 \Um$ 处乙波的平衡位置振动刚好传到 P 处，且此时两列波的振动都向 $y$ 轴正方向运动，根据叠加原理可知，$t = 1.0 \Us$ 时，P 向 $y$ 轴正方向运动，故 C 正确，D 错误。

故选 BC。
}

\item
%% number: 10
%% paperName: 2024年山东高考第 10 题 4 分
%% typeId: 2
%% type: 多选题
%% chapter: 电场
%% point: 带电粒子的直线运动
%% method: 无
%% score: 4
%% degree: 450
%% duplicateId: 0
%% body:
如图所示，带电量为 $+q$ 的小球被绝缘棒固定在 O 点，右侧有固定在水平面上、倾角为 $30^{\circ}$ 的光滑绝缘斜面。质量为 $m$、带电量为 $+q$ 的小滑块从斜面上 A 点由静止释放，滑到与小球等高的 B 点时加速度为零，滑到 C 点时速度为零。已知 AC 间的距离为 $s$，重力加速度大小为 $g$，静电力常量为 $k$，下列说法正确的是\xzanswer{AD}

\onepicture[w=5.50cm, align=right]{figs/10.png}

\fourchoices[answer=AD]
{OB 的距离 $l = \sqrt{\frac{\sqrt{3}kq^{2}}{mg}}$}
{OB 的距离 $l = \sqrt{\frac{\sqrt{3}kq^{2}}{3mg}}$}
{从 A 到 C，静电力对小滑块做功 $W = - mgs$}
{AC 之间的电势差 $U_{AC} = - \frac{mgs}{2q}$}

%% answer: AD
%% memo:
\memoanswer{
AB．由题意知小滑块在 B 点处的加速度为零，则根据受力分析有沿斜面方向

$mg\sin30^{\circ} = \frac{kq^{2}}{l^{2}}\cos30^{\circ}$

解得 $l = \sqrt{\frac{\sqrt{3}kq^{2}}{mg}}$

A 正确，B 错误；

C．因为滑到 C 点时速度为零，小滑块从 A 到 C 的过程，静电力对小滑块做的功为 $W$，根据动能定理有

$W + mgs\sin30^{\circ} = 0$

解得 $W = - \frac{mgs}{2}$

故 C 错误；

D．根据电势差与电场强度的关系可知 AC 之间的电势差

$U_{AC} = \frac{W}{q} = - \frac{mgs}{2q}$

故 D 正确。

故选 AD。
}

\item
%% number: 11
%% paperName: 2024年山东高考第 11 题 4 分
%% typeId: 2
%% type: 多选题
%% chapter: 电磁感应
%% point: 电磁感应能量分析
%% method: 无
%% score: 4
%% degree: 550
%% duplicateId: 0
%% body:
如图所示，两条相同的半圆弧形光滑金属导轨固定在水平桌面上，其所在平面竖直且平行，导轨最高点到水平桌面的距离等于半径，最低点的连线 $OO'$ 与导轨所在竖直面垂直。空间充满竖直向下的匀强磁场（图中未画出），导轨左端由导线连接。现将具有一定质量和电阻的金属棒 MN 平行 $OO'$ 放置在导轨图示位置，由静止释放。MN 运动过程中始终平行于 $OO'$ 且与两导轨接触良好，不考虑自感影响，下列说法正确的是\xzanswer{ABD}

\onepicture[w=7.00cm, align=right]{figs/11.png}

\fourchoices[answer=ABD]
{MN 最终一定静止于 $OO'$ 位置}
{MN 运动过程中安培力始终做负功}
{从释放到第一次到达 $OO'$ 位置过程中，MN 的速率一直在增大}
{从释放到第一次到达 $OO'$ 位置过程中，MN 中电流方向由 M 到 N}

%% answer: ABD
%% memo:
\memoanswer{
A．由于金属棒 MN 运动过程切割磁感线产生感应电动势，回路有感应电流，产生焦耳热，金属棒 MN 的机械能不断减小，由于金属导轨光滑，所以经过多次往返运动，MN 最终一定静止于 $OO'$ 位置，故 A 正确；

B．当金属棒 MN 向右运动，根据右手定则可知，MN 中电流方向由 M 到 N，根据左手定则，可知金属棒 MN 受到的安培力水平向左，则安培力做负功；当金属棒 MN 向左运动，根据右手定则可知，MN 中电流方向由 N 到 M，根据左手定则，可知金属棒 MN 受到的安培力水平向右，则安培力做负功；可知 MN 运动过程中安培力始终做负功，故 B 正确；

C．金属棒 MN 从释放到第一次到达 $OO'$ 位置过程中，由于在 $OO'$ 位置重力沿切线方向的分力为 0，可知在到达 $OO'$ 位置之前的位置，重力沿切线方向的分力已经小于安培力沿切线方向的分力，金属棒 MN 已经做减速运动，故 C 错误；

D．从释放到第一次到达 $OO'$ 位置过程中，根据右手定则可知，MN 中电流方向由 M 到 N，故 D 正确。

故选 ABD。
}

\item
%% number: 12
%% paperName: 2024年山东高考第 12 题 4 分
%% typeId: 2
%% type: 多选题
%% chapter: 曲线运动
%% point: 斜抛运动
%% method: 无
%% score: 4
%% degree: 350
%% duplicateId: 0
%% body:
如图所示，工程队向峡谷对岸平台抛射重物，初速度 $v_{0}$ 大小为 $20 \Ums$，与水平方向的夹角为 $30^{\circ}$，抛出点 P 和落点 Q 的连线与水平方向夹角为 $30^{\circ}$，重力加速度大小取 $10 \Umsq$，忽略空气阻力。重物在此运动过程中，下列说法正确的是\xzanswer{BD}

\onepicture[w=5.00cm, align=right]{figs/12.png}

\fourchoices[answer=BD]
{运动时间为 $2\sqrt{3} \Us$}
{落地速度与水平方向夹角为 $60^{\circ}$}
{重物离 PQ 连线的最远距离为 $10 \Um$}
{轨迹最高点与落点的高度差为 $45 \Um$}

%% answer: BD
%% memo:
\memoanswer{
AC．将初速度分解为沿 PQ 方向分速度 $v_{1}$ 和垂直 PQ 分速度 $v_{2}$，则有 $v_{1}=v_{0}\cos60^{\circ}=10 \Ums$，$v_{2}=v_{0}\sin60^{\circ}=10\sqrt{3} \Ums$

将重力加速度分解为沿 PQ 方向分速度 $a_{1}$ 和垂直 PQ 分速度 $a_{2}$，则有

$a_{1}=g\sin30^{\circ}=5 \Umsq$，$a_{2}=g\cos30^{\circ}=5\sqrt{3} \Umsq$

垂直 PQ 方向根据对称性可得重物运动时间为

$t=2\frac{v_{2}}{a_{2}}=4 \Us$

重物离 PQ 连线的最远距离为

$d_{\max}=\frac{v_{2}^{2}}{2a_{2}}=10\sqrt{3} \Um$

故 AC 错误；

B．重物落地时竖直分速度大小为

$v_{y}=-v_{0}\sin30^{\circ}+gt=30 \Ums$

则落地速度与水平方向夹角正切值为

$\tan\theta=\frac{v_{y}}{v_{x}}=\frac{v_{y}}{v_{0}\cos30^{\circ}}=\sqrt{3}$

可得 $\theta=60^{\circ}$

故 B 正确；

D．从抛出到最高点所用时间为

$t_{1}=\frac{v_{0}\sin30^{\circ}}{g}=1 \Us$

则从最高点到落地所用时间为

$t_{2}=t-t_{1}=3 \Us$

轨迹最高点与落点的高度差为

$h=\frac{1}{2}gt_{2}^{2}=45 \Um$

故 D 正确。

故选 BD。
}

\item
%% number: 13
%% paperName: 2024年山东高考第 13 题 6 分
%% typeId: 4
%% type: 实验题
%% chapter: 运动和力
%% point: 验证动量守恒实验
%% method: 无
%% score: 6
%% degree: 650
%% duplicateId: 0
%% body:
在第四次“天宫课堂”中，航天员演示了动量守恒实验。受此启发，某同学使用如图甲所示的装置进行了碰撞实验，气垫导轨两端分别安装 a、b 两个位移传感器，a 测量滑块 A 与它的距离 $x_{A}$，b 测量滑块 B 与它的距离 $x_{B}$。部分实验步骤如下：

①测量两个滑块的质量，分别为 $200.0 \Ug$ 和 $400.0 \Ug$；

②接通气源，调整气垫导轨水平；

③拨动两滑块，使 A、B 均向右运动；

④导出传感器记录的数据，绘制 $x_{A}$、$x_{B}$ 随时间变化的图像，分别如图乙、图丙所示。

回答以下问题：
\begin{enumerate}
	\item
	从图像可知两滑块在 $t = $\tkanswer[1.5cm]{1.0}$\Us$ 时发生碰撞；
	\item
	滑块 B 碰撞前的速度大小 $v = $\tkanswer[1.5cm]{0.20}$\Ums$（保留 2 位有效数字）；
	\item
	通过分析，得出质量为 $200.0 \Ug$ 的滑块是\tkanswer[1cm]{B}（填“A”或“B”）。
\end{enumerate}
\onepicture[w=9.00cm, align=right]{figs/13.png}
%% answer:
\jdanswer{
\begin{enumerate}
	\item
	$1.0$

	\item
	$0.20$

	\item
	B

\end{enumerate}
}
%% memo:
\memoanswer{
（1）由 $x-t$ 图像的斜率表示速度可知两滑块的速度在 $t = 1.0 \Us$ 时发生突变，即这个时候发生了碰撞；

（2）根据 $x-t$ 图像斜率的绝对值表示速度大小可知碰撞前瞬间 B 的速度大小为

$v=\left|\frac{90-110}{1.0}\right|\Ucms=0.20 \Ums$

（3）由题图乙知，碰撞前 A 的速度大小 $v_{A}=0.50 \Ums$，碰撞后 A 的速度大小约为 $v_{A}'=0.36 \Ums$，由题图丙可知，碰撞后 B 的速度大小为 $v_{B}'=0.5 \Ums$，A 和 B 碰撞过程动量守恒，则有

$m_{A}v_{A}+m_{B}v=m_{A}v_{A}'+m_{B}v_{B}'$

代入数据解得 $\frac{m_{A}}{m_{B}} \approx 2$

所以质量为 $200.0 \Ug$ 的滑块是 B。
}

\item
%% number: 14
%% paperName: 2024年山东高考第 14 题 8 分
%% typeId: 4
%% type: 实验题
%% chapter: 电路
%% point: 测量电阻率
%% method: 无
%% score: 8
%% degree: 550
%% duplicateId: 0
%% body:
某学习小组对两种型号铅笔芯的电阻率进行测量。实验器材如下：

学生电源（输出电压 $0 \sim 16 \UV$）；

滑动变阻器（最大阻值 $10 \UO$，额定电流 $2 \UA$）；

电压表 V（量程 $3 \UV$，内阻未知）；

电流表 A（量程 $3 \UA$，内阻未知）；

待测铅笔芯 $R$（X 型号、Y 型号）；

游标卡尺，螺旋测微器，开关 S，单刀双掷开关 K，导线若干。

回答以下问题：
\begin{enumerate}
	\item
	使用螺旋测微器测量铅笔芯直径，某次测量结果如图甲所示，该读数为\tkanswer[1.5cm]{2.450}$\Umm$；
	\item
	把待测铅笔芯接入图乙所示电路，闭合开关 S 后，将滑动变阻器滑片由最右端向左调节到合适位置，将单刀双掷开关 K 分别掷到 1、2 端，观察到电压表示数变化比电流表示数变化更明显，则测量铅笔芯电阻时应将 K 掷到\tkanswer[1cm]{1}（填“1”或“2”）端；
	\item
	正确连接电路，得到 Y 型号铅笔芯 $I - U$ 图像如图丙所示，求得电阻 $R_{Y} = $\tkanswer[1.5cm]{1.91}$\UO$（保留 3 位有效数字）；采用同样方法得到 X 型号铅笔芯的电阻为 $1.70 \UO$；
	\item
	使用游标卡尺测得 X、Y 型号铅笔芯的长度分别为 $40.68 \Umm$、$60.78 \Umm$，使用螺旋测微器测得 X、Y 型号铅笔芯直径近似相等，则 X 型号铅笔芯的电阻率\tkanswer[1cm]{大于}（填“大于”或“小于”）Y 型号铅笔芯的电阻率。
\end{enumerate}
\onepicture[w=9.00cm, align=right]{figs/14a.png}

\onepicture[w=5.50cm, align=right]{figs/14b.png}
%% answer:
\jdanswer{
\begin{enumerate}
	\item
	$2.450$

	\item
	$1$

	\item
	$1.91$

	\item
	大于

\end{enumerate}
}
%% memo:
\memoanswer{
（1）根据螺旋测微器的读数规则可知，其读数为

$d = 2 \Umm + 0.01 \times 45.0 \Umm = 2.450 \Umm$

（2）由于电压表示数变化更明显，说明电流表分压较多，因此电流表应采用外接法，即测量铅笔芯电阻时应将 K 掷到 1 端；

（3）根据图丙的 $I - U$ 图像，结合欧姆定律有

$R_{Y}=\frac{2.50 \UV}{1.31 \UA}=1.91 \UO$

（4）根据电阻定律 $R = \rho \frac{l}{S}$

可得 $\rho=\frac{RS}{l}$

两种材料的横截面积近似相等，分别代入数据可知 $\rho_{X} > \rho_{Y}$。
}

\item
%% number: 15
%% paperName: 2024年山东高考第 15 题 8 分
%% typeId: 6
%% type: 计算题
%% chapter: 光学
%% point: 全反射
%% method: 无
%% score: 8
%% degree: 550
%% duplicateId: 0
%% body:
某光学组件横截面如图所示，半圆形玻璃砖圆心为 O 点，半径为 $R$；直角三棱镜 FG 边的延长线过 O 点，EG 边平行于 AB 边且长度等于 $R$，$\angle FEG = 30^{\circ}$。横截面所在平面内，单色光线以 $\theta$ 角入射到 EF 边发生折射，折射光线垂直 EG 边射出。已知玻璃砖和三棱镜对该单色光的折射率均为 1.5。
\begin{enumerate}
	\item
	求 $\sin\theta$；
	\item
	以 $\theta$ 角入射的单色光线，若第一次到达半圆弧 AMB 可以发生全反射，求光线在 EF 上入射点 D（图中未标出）到 E 点距离的范围。
\end{enumerate}
\onepicture[w=5.20cm, align=right]{figs/15.png}
%% answer:
\jdanswer{
\begin{enumerate}
	\item
	$\sin\theta = 0.75$

	\item
	$0 < x \le \frac{2\sqrt{3}}{9}R$

\end{enumerate}
}
%% memo:
\memoanswer{
（1）由题意设光在三棱镜中的折射角为 $\alpha$，则根据折射定律有

$n=\frac{\sin\theta}{\sin\alpha}$

由于折射光线垂直 EG 边射出，根据几何关系可知

$\alpha = \angle FEG = 30^{\circ}$

代入数据解得 $\sin\theta = 0.75$

（2）根据题意作出单色光第一次到达半圆弧 AMB 恰好发生全反射的光路图如图。

\onepicture[w=5.00cm]{figs/15_an.jpg}

则根据几何关系可知 FE 上从 P 点到 E 点以 $\theta$ 角入射的单色光线第一次到达半圆弧 AMB 都可以发生全反射，根据全反射临界角公式有

$\sin C = \frac{1}{n}$

设 P 点到 FG 的距离为 $l$，则根据几何关系有 $l = R \sin C$

又因为 $x_{PE} = \frac{R - l}{\cos30^{\circ}}$

联立解得 $x_{PE} = \frac{2\sqrt{3}}{9}R$

所以光线在 EF 上的入射点 D 到 E 点的距离范围为 $0 < x \le \frac{2\sqrt{3}}{9}R$。
}

\item
%% number: 16
%% paperName: 2024年山东高考第 16 题 8 分
%% typeId: 6
%% type: 计算题
%% chapter: 气体性质
%% point: 玻意耳定律
%% method: 无
%% score: 8
%% degree: 550
%% duplicateId: 0
%% body:
图甲为战国时期青铜汲酒器，根据其原理制作了由中空圆柱形长柄和储液罐组成的汲液器，如图乙所示。长柄顶部封闭，横截面积 $S_{1} = 1.0 \Ucmq$，长度 $H = 100.0 \Ucm$，侧壁有一小孔 A。储液罐的横截面积 $S_{2} = 90.0 \Ucmq$，高度 $h = 20.0 \Ucm$，罐底有一小孔 B。汲液时，将汲液器竖直浸入液体，液体从孔 B 进入，空气由孔 A 排出；当内外液面相平时，长柄浸入液面部分的长度为 $x$；堵住孔 A，缓慢地将汲液器竖直提出液面，储液罐内刚好储满液体。已知液体密度 $\rho = 1.0 \times 10^{3} \Ukgmc$，重力加速度大小 $g = 10 \Umsq$，大气压 $p_{0} = 1.0 \times 10^{5} \UPa$。整个过程温度保持不变，空气可视为理想气体，忽略器壁厚度。
\begin{enumerate}
	\item
	求 $x$；
	\item
	松开孔 A，从外界进入压强为 $p_{0}$、体积为 $V$ 的空气，使满储液罐中液体缓缓流出，堵住孔 A，稳定后罐中恰好剩余一半的液体，求 $V$。
\end{enumerate}
\onepicture[w=5.50cm, align=right]{figs/16.png}
%% answer:
\jdanswer{
\begin{enumerate}
	\item
	$x = 2 \Ucm$

	\item
	$V = 8.92 \times 10^{-4} \Umc$

\end{enumerate}
}
%% memo:
\memoanswer{
（1）由题意可知缓慢地将汲液器竖直提出液面过程，气体发生等温变化，所以有

$p_{1}(H-x)S_{1}=p_{2}HS_{1}$

又因为 $p_{1}=p_{0}$，$p_{2}+\rho gh=p_{0}$

代入数据联立解得 $x = 2 \Ucm$

（2）当外界气体进入后，以所有气体为研究对象有

$p_{0}V+p_{2}HS_{1}=p_{3}\left(HS_{1}+\frac{h}{2}S_{2}\right)$

又因为 $p_{3} + \rho g \cdot \frac{h}{2} = p_{0}$

代入数据联立解得 $V = 8.92 \times 10^{-4} \Umc$。
}

\item
%% number: 17
%% paperName: 2024年山东高考第 17 题 14 分
%% typeId: 6
%% type: 计算题
%% chapter: 运动和力
%% point: 动量和能量
%% method: 无
%% score: 14
%% degree: 450
%% duplicateId: 0
%% body:
如图甲所示，质量为 $M$ 的轨道静止在光滑水平面上，轨道水平部分的上表面粗糙，竖直半圆形部分的表面光滑，两部分在 P 点平滑连接，Q 为轨道的最高点。质量为 $m$ 的小物块静置在轨道水平部分上，与水平轨道间的动摩擦因数为 $\mu$，最大静摩擦力等于滑动摩擦力。已知轨道半圆形部分的半径 $R = 0.4 \Um$，重力加速度大小 $g = 10 \Umsq$。
\begin{enumerate}
	\item
	若轨道固定，小物块以一定的初速度沿轨道运动到 Q 点时，受到轨道的弹力大小等于 $3mg$，求小物块在 Q 点的速度大小 $v$；
	\item
	若轨道不固定，给轨道施加水平向左的推力 $F$，小物块处在轨道水平部分时，轨道加速度 $a$ 与 $F$ 对应关系如图乙所示。（i）求 $\mu$ 和 $m$；（ii）初始时，小物块静置在轨道最左端，给轨道施加水平向左的推力 $F = 8 \UN$，当小物块到 P 点时撤去 $F$，小物块从 Q 点离开轨道时相对地的速度大小为 $7 \Ums$。求轨道水平部分的长度 $L$。
\end{enumerate}
\onepicture[w=9.00cm, align=right]{figs/17.png}
%% answer:
\jdanswer{
\begin{enumerate}
	\item
	$v = 4 \Ums$

	\item
	（i）$m = 1 \Ukg$，$\mu = 0.2$；（ii）$L = 4.5 \Um$

\end{enumerate}
}
%% memo:
\memoanswer{
（1）根据题意可知小物块在 Q 点由合力提供向心力有

$mg+3mg=m\frac{v^{2}}{R}$

代入数据解得 $v = 4 \Ums$

（2）（i）根据题意可知当 $F \le 4 \UN$ 时，小物块与轨道是一起向左加速，根据牛顿第二定律可知 $F = (M + m)a$

根据图乙有 $k=\frac{1}{M+m}=0.5\Ukg^{-1}$

当外力 $F > 4 \UN$ 时，轨道与小物块有相对滑动，则对轨道有 $F - \mu mg = Ma$

结合题图乙有 $a=\frac{1}{M}F-\frac{\mu mg}{M}$

可知 $k=\frac{1}{M}=1\Ukg^{-1}$

截距 $b = -\frac{\mu mg}{M} = -2 \Umsq$

联立以上各式可得 $M = 1 \Ukg$，$m = 1 \Ukg$，$\mu = 0.2$

（ii）由图乙可知，当 $F = 8 \UN$ 时，轨道的加速度为 $6 \Umsq$，小物块的加速度为 $a_{2} = \mu g = 2 \Umsq$

当小物块运动到 P 点时，经过 $t_{0}$ 时间，则轨道有 $v_{1} = a_{1}t_{0}$

小物块有 $v_{2} = a_{2}t_{0}$

在这个过程中系统机械能守恒有

$\frac{1}{2}Mv_{1}^{2} + \frac{1}{2}mv_{2}^{2} = \frac{1}{2}Mv_{3}^{2} + \frac{1}{2}mv_{4}^{2} + 2mgR$

水平方向动量守恒，以水平向左为正方向，则有

$Mv_{1} + mv_{2} = Mv_{3} + mv_{4}$

联立解得 $t_{0} = 1.5 \Us$

根据运动学公式有 $L = \frac{1}{2}a_{1}t_{0}^{2} - \frac{1}{2}a_{2}t_{0}^{2}$

代入数据解得 $L = 4.5 \Um$。
}

\item
%% number: 18
%% paperName: 2024年山东高考第 18 题 16 分
%% typeId: 6
%% type: 计算题
%% chapter: 磁场
%% point: 带电粒子在组合场中的运动
%% method: 无
%% score: 16
%% degree: 250
%% duplicateId: 0
%% body:
如图所示，在 $Oxy$ 坐标系 $x > 0$，$y > 0$ 区域内充满垂直纸面向里，磁感应强度大小为 $B$ 的匀强磁场。磁场中放置一长度为 $L$ 的挡板，其两端分别位于 $x$、$y$ 轴上 M、N 两点，$\angle OMN = 60^{\circ}$，挡板上有一小孔 K 位于 MN 中点。$\triangle OMN$ 之外的第一象限区域存在恒定匀强电场。位于 $y$ 轴左侧的粒子发生器在 $0 < y < \frac{\sqrt{3}}{2}L$ 的范围内可以产生质量为 $m$，电荷量为 $+q$ 的无初速度的粒子。粒子发生器与 $y$ 轴之间存在水平向右的匀强加速电场，加速电压大小可调，粒子经此电场加速后进入磁场，挡板厚度不计，粒子可沿任意角度穿过小孔，碰撞挡板的粒子不予考虑，不计粒子重力及粒子间相互作用力。
\begin{enumerate}
	\item
	求使粒子垂直挡板射入小孔 K 的加速电压 $U_{0}$；
	\item
	调整加速电压，当粒子以最小的速度从小孔 K 射出后恰好做匀速直线运动，求第一象限中电场强度的大小和方向；
	\item
	当加速电压为 $\frac{qB^{2}L^{2}}{24m}$ 时，求粒子从小孔 K 射出后，运动过程中距离 $y$ 轴最近位置的坐标。
\end{enumerate}
\onepicture[w=5.20cm, align=right]{figs/18.png}
%% answer:
\jdanswer{
\begin{enumerate}
	\item
	$U_{0} = \frac{qB^{2}L^{2}}{8m}$

	\item
	$E = \frac{qB^{2}L}{4m}$，方向沿 $x$ 轴正方向

	\item
	$d_{m} = \frac{(3 - \sqrt{3})L}{12}$

\end{enumerate}
}
%% memo:
\memoanswer{
（1）根据题意，作出粒子垂直挡板射入小孔 K 的运动轨迹如图所示。

\onepicture[w=5.00cm]{figs/18_an1.jpg}

根据几何关系可知粒子在磁场中做圆周运动的轨迹半径为

$r = x_{SK} = \frac{L}{2}$

在 $\triangle OMN$ 区域根据洛伦兹力提供向心力有

$qvB = m\frac{v^{2}}{r}$

在匀强加速电场中由动能定理有

$U_{0}q = \frac{1}{2}mv^{2}$

联立解得 $U_{0} = \frac{qB^{2}L^{2}}{8m}$

（2）根据题意，当轨迹半径最小时，粒子速度最小，则作出粒子以最小的速度从小孔 K 射出的运动轨迹如图所示。

\onepicture[w=5.00cm]{figs/18_an2.jpg}

根据几何关系可知粒子在磁场中做圆周运动的轨迹半径为

$r' = x_{NK}\cos60^{\circ} = \frac{L}{4}$

在 $\triangle OMN$ 区域根据洛伦兹力提供向心力有

$qv'B = m\frac{v'^{2}}{r'}$

粒子从小孔 K 射出后恰好做匀速直线运动，由左手定则可知粒子经过小孔 K 后受到的洛伦兹力沿 $x$ 轴负方向，则粒子经过小孔 K 后受到的电场力沿 $x$ 轴正方向，又粒子带正电，则 $\triangle OMN$ 之外第一象限区域电场强度的方向沿 $x$ 轴正方向，大小满足

$qv'B = Eq$

联立可得 $E = \frac{qB^{2}L}{4m}$

（3）在匀强加速电场中由动能定理有

$Uq = \frac{1}{2}mv''^{2}$

可得 $v'' = \frac{\sqrt{3}qBL}{6m}$

在 $\triangle OMN$ 区域根据洛伦兹力提供向心力有

$qv''B = m\frac{v''^{2}}{r''}$

可得粒子在 $\triangle OMN$ 区域运动的轨迹半径

$r'' = \frac{\sqrt{3}}{6}L$

作出从小孔 K 射出的粒子的运动轨迹如图所示。

\onepicture[w=5.00cm]{figs/18_an3.jpg}

粒子从 K 射出时，$v''$ 越偏向 $y$ 轴，离 $y$ 轴越近，由几何关系有

$\sin\theta = \frac{\frac{L}{4}}{r''} = \frac{\sqrt{3}}{2}$

则有 $\theta = 60^{\circ}$

由配速法将运动分解为 $y$ 轴方向的匀速直线运动和沿 $x$ 方向的匀速圆周运动，其中

$v''_{y} = v''\sin\theta = \frac{qBL}{4m}$

$v''_{x} = v''\cos\theta = \frac{qBL}{4\sqrt{3}m}$

匀速圆周运动的半径为

$r_{y} = \frac{mv''_{x}}{qB} = \frac{\sqrt{3}}{12}L$

故最小距离为

$d_{m} = \frac{L}{4} - r_{y} = \frac{(3 - \sqrt{3})L}{12}$。
}

\end{enumerate}

\end{document}
